%% ma: L11-B15-0087
%% lop: 11
%% chuong: 5
%% bai: 15
%% dang: 11.15.7
%% loai: TLN
%% muc: TH
%% dap_an: "16"
%% nguon: "tai-lieu-hoc-tap-mon-toan-lop-11-theo-chuong-trinh-giao-duc-pho-thong-2018.pdf (Lớp Toán Thầy Duy, Phạm Lê Duy, THPT Chu Văn An), Chương 3 Bài 1, Phần 3, Câu 9 (tr.343)"
%% kien_thuc: "Lớp 11 Bài 15 (giới hạn của dãy số)"
%% trang_thai: cho-duyet
%% ly_do: "Dạng mới đề xuất (11.15.x), chờ giáo viên duyệt dạng; tài liệu gốc không có lời giải mẫu, đáp án do kho tự giải và kiểm bằng sympy."
\begin{ex}
Giá trị của tổng \[T=1+\dfrac1{\sqrt2}+\dfrac12+\dfrac1{2\sqrt2}+\cdots+\left(\dfrac1{\sqrt2}\right)^n+\cdots=a+\sqrt b\] với $a;b$ là các số tự nhiên. Tính giá trị $S=(a+b)^2$.
\shortans{16}
\loigiai{\sloppy Cấp số nhân lùi vô hạn có $u_1=1$, $q=\dfrac1{\sqrt2}$: $T=\dfrac1{1-\frac1{\sqrt2}}=\dfrac{\sqrt2}{\sqrt2-1}=2+\sqrt2$. Vậy $a=2$, $b=2$ và $S=(2+2)^2=16$.}
\end{ex}
